\begin{textsample}{POS Dim 2 – 2013\_c5 – Score -1.00 – 2013\_c5/enade\_2013\_zoot\_9cdde507ba.txt}  \label{ex:f2_pos_366}
CAPÍTULO 5 DISTRIBUIÇÃO DOS CONCEITOS 5.1 PANORAMA NACIONAL DA DISTRIBUIÇÃO DOS CONCEITOS A Tabela 5.1 apresenta a quantidade e distribuição de cursos de Zootecnia participantes do ENADE/2013, por faixa de conceito e Grande Região.
A diferença entre os cursos tabulados neste capítulo e no capítulo 2 são os cursos sem conceito, em princípio, aqueles sem alunos concluintes que participassem da prova ou que tivessem tão somente um aluno concluinte13.
Observando -se os dados da Tabela 5.1, nota -se que, dos 82 cursos participantes, 30 ( 36,6 \% ) classificaram -se com conceito 3, o valor modal.
Este foi também o conceito modal em todas as regiões do Brasil ( entre 30,0 \% na região Norte e 50,0 \% na região Centro-Oeste ), exceto na região Sudeste, onde a moda foi o conceito 4 ( 60,0 \% ).
Vale ressaltar que na região Sul além do conceito 3, já citado, ocorreu outro valor modal correspondente ao conceito 4.
Portanto, os conceitos nessa região caracterizaram -se por a moda corresponder a dois conceitos contíguos.
O conceito 4 foi o segundo mais frequente em nível nacional ( 34,1 \%, correspondendo a 28 cursos ), e os conceitos 2 e 5 ficaram, ambos em terceiro ( 9,8 \%, cada, correspondendo a oito cursos, cada ).
Houve, ainda, seis cursos que receberam conceito 1.
Dos 82 cursos de Zootecnia, apenas dois ficaram sem conceito ( SC ).
A região Norte participou com dez cursos ou 12,2 \% do total nacional.
Nesta região, o conceito 3 foi atribuído a três cursos dos dez participantes, o que equivale a 30,0 \% do total regional.
Apenas um curso ficou sem conceito ( 10,0 \% ).
Os demais foram igualmente 13 Estes cursos com somente um concluinte não foram considerados no capítulo 2. distribuídos e avaliados com conceitos 1, 2 e 4 ( dois cursos, cada, correspondendo a 20,0 \%, cada ).
Nenhum curso foi alocado no conceito 5.
A região Nordeste participou com 24 cursos ou 29,3 \% do total nacional.
Nesta região, 11 cursos, 45,8 \% em termos regionais, obtiveram conceito 3, o conceito modal para a região.
Os conceitos 2 e 4 foram atribuídos a, respectivamente, três e cinco cursos ( 12,5 \% e 20,8 \% ).
Já os conceitos 1 e 5 foram atribuídos a, respectivamente, quatro e um cursos ( 16,7 \% e 4,2 \% ).
Nessa região, nenhum dos cursos ficou sem conceito ( SC ).
Como já comentado, dos 20 cursos participantes da região Sudeste, 12 ( 60,0 \% ) obtiveram conceito 4, o conceito modal.
Nenhum curso ficou sem conceito ( SC ) ou obteve conceitos 1 ou 2.
O conceito 3 foi atribuído a três cursos ( 15,0 \% ) e cinco outros ( 25,0 \% ) receberam o conceito 5.
Na região Sul houve 14 cursos participantes, com a predominância nos conceitos 3 e 4 correspondentes a 42,9 \%, cada, caracterizando, portanto, conceito uma moda na região, distribuída entre dois conceitos contíguos, como já comentado anteriormente.
O conceito 5 foi atribuído apenas a dois cursos ( 14,3 \% ).
Nessa região, nenhum curso ficou sem conceito ( SC ) ou recebeu os conceitos 1 ou 2.
Como já comentado, o conceito modal na região Centro-Oeste foi 3, sete dos 14 cursos participantes receberam este conceito ( 50,0 \% dos cursos ).
Nesta região somente um curso ficou sem conceito ( 7,1 \% ).
Já, três cursos foram alocados no conceito 2 e outros três no conceito 4, correspondendo juntos a 42,8 \% dos cursos da região.
Os conceitos 1 e 5 não foram alocados a nenhum curso.
O Gráfico 5.1 apresenta a distribuição do Conceito ENADE segundo Grande Região : poligonais mais à esquerda representam Grandes Regiões com pior distribuição, e poligonais mais à direita, Grandes Regiões com melhores conceitos.
As regiões Norte ( linha azul ), Nordeste ( linha verde ) e Centro-Oeste ( linha preta ) apresentam os cursos com valores mais baixos, ao passo que as regiões Sul ( linha roxa ) e Sudeste ( linha vermelha ) apresentam as distribuições com valores maiores. 5.2 CONCEITOS POR CATEGORIA ADMINISTRATIVA E POR GRANDE REGIÃO A Tabela 5.2 apresenta a distribuição dos cursos participantes do ENADE/2013 de Zootecnia, por Categoria Administrativa, de acordo com os conceitos por eles alcançados, segundo as Grandes Regiões brasileiras.
Dos 82 cursos participantes, 68 ( 82,9 \% ) eram ministrados em Instituições Públicas e 14 ( 17,1 \% ), em Instituições Privadas.
De acordo com as informações da Tabela 5.2, em termos nacionais, entre Instituições Públicas e Privadas : das oito IES que receberam o conceito 5, todas eram Públicas.
Dos 68 cursos participantes de IES Públicas, o conceito 4 foi o valor modal, atribuído a 27 cursos ( 39,7 \% ).
Entre os demais cursos de Instituições Públicas participantes, três obtiveram conceito 1 ( 4,4 \% da categoria ), cinco receberam conceito 2 ( 7,4 \% ) e 24 foram alocados no conceito 3 ( 35,3 \% ).
Como já comentado, os oito restantes receberam o conceito 5 ( 11,8 \% ).
Nesta categoria, apenas um dos cursos ficou sem conceito ( SC ).
Na rede privada, o conceito modal também foi 3, com seis cursos dos 14 da categoria.
Entre os demais cursos participantes, três cursos receberam conceito 1 e outros três, conceito 2.
O conceito 4 foi atribuído a somente um curso e o conceito 5, como já comentado, a nenhum curso.
Nesta Categoria Administrativa, apenas um dos cursos ficou sem conceito ( SC ).
Na análise por região, observa -se que, na região Norte, as Instituições Privadas participaram com quatro cursos ( 40,0 \% do total regional ), proporção mais elevada do que nas demais regiões brasileiras.
Nessa região, nenhum curso ficou com conceitos 3, 4 ou 5.
O conceito modal para as Instituições Privadas na região foi 2, com dois cursos, correspondendo a 50,0 \%.
O conceito 1 apresentou somente um curso, correspondendo a 25,0 \%, e apenas um curso ficou sem conceito ( SC ).
As Instituições Públicas participaram com seis cursos ( 60,0 \% do total regional ), dos quais nenhum ficou sem conceito ( SC ) ou recebeu os conceitos 2 ou 5 nesta categoria e região.
O conceito modal para esta Categoria Administrativa, na região Norte, foi 3, com três cursos, correspondendo a 50,0 \%.
Os conceitos 1 e 4 participaram com um e dois cursos ( 16,7 \% e 33,3 \% ), respectivamente.
Na região Nordeste, a rede privada concentrou dois dos 24 cursos participantes, equivalentes a 8,3 \% do total da região, a segunda menor proporção desta categoria entre todas as regiões.
Essa Categoria Administrativa na região concentrou todos os dois cursos com conceito 2, correspondendo à totalidade de cursos oferecidos por IES Privadas no Nordeste.
As Instituições Públicas dessa região participaram com 22 cursos ( 91,7 \% ), dos quais 11 obtiveram conceitos 3, o valor modal.
Os demais foram avaliados com os conceitos 1 ( dois cursos ), 2 ( três cursos ), 4 ( cinco cursos ) e 5 ( um curso ).
Nenhum curso ficou sem conceito ( SC ) nesta combinação de região e categoria.
Na região Sudeste, a proporção de cursos da rede privada foi 20,0 \%, correspondendo a quatro dos 20 cursos participantes.
Nesta categoria, na região Sudeste, o conceito modal foi 3 ( três cursos ).
Nesta combinação de Categoria Administrativa e Grande Região, apenas um curso recebeu conceito 4.
Nenhum curso ficou sem conceito ( SC ) ou recebeu os demais conceitos nessa categoria e região.
Entre os 16 cursos em Instituições Públicas na região Sudeste, a categoria modal foi a que recebeu conceito 4 ( 11 cursos ), enquanto cinco cursos obtiveram conceito 5.
E também, na região Sudeste, nesta categoria, nenhum curso ficou sem conceito ( SC ) ou recebeu os demais conceitos.
As Instituições Privadas concentraram somente um dos 14 cursos participantes da região Sul, 7,1 \% do total regional, a menor proporção desta categoria entre todas as regiões.
Esse único curso foi alocado ao conceito 3, correspondendo, portanto, à totalidade de cursos nessa categoria e região.
As Instituições Públicas na região Sul participaram com 13 cursos ( 92,9 \% ).
O conceito modal foi 4, com seis cursos.
Já cinco cursos obtiveram conceito 3, enquanto apenas dois cursos receberam conceito 5.
Nenhum curso ficou sem conceito ( SC ) ou recebeu os conceitos 1 ou 2.
Na região Centro-Oeste, apenas três dos 14 cursos participantes eram de Instituições Privadas ( 21,4 \% em termos regionais ).
Destes, dois concentraram -se no conceito 3, conceito modal, e somente um curso foi alocado no conceito 2.
Nessa categoria e região nenhum curso ficou sem conceito ( SC ) ou recebeu os demais conceitos restantes.
Dos 11 cursos de Instituições Públicas, cinco foram avaliados no conceito 3, conceito modal.
Os demais receberam conceito 2 ( dois cursos ) e conceito 4 ( três cursos ).
Nesta região, somente um curso ficou sem conceito ( SC ), enquanto nenhum curso de IES Pública recebeu os conceitos 1 ou 5.
O Gráfico 5.2 apresenta a distribuição do Conceito ENADE segundo a Categoria Administrativa da IES.
Os cursos em IES Públicas ( linha azul ) apresentam uma poligonal mais à direita do que os em IES Privadas ( linha verde ) e, consequentemente, uma distribuição de Conceitos ENADE com valores maiores. 5.3 CONCEITOS POR ORGANIZAÇÃO ACADÊMICA E POR GRANDE REGIÃO Na Tabela 5.3 encontra -se a distribuição dos conceitos atribuídos aos cursos participantes do ENADE/2013 na Área de Zootecnia, por Organização Acadêmica, segundo as Grandes Regiões brasileiras.
Dos 82 cursos participantes, 71 eram oferecidos em Universidades, dez em Faculdades e um em Centros Universitários.
Esta distribuição corresponde a, respectivamente, 86,6 \%, 12,2 \% e 1,2 \% dos cursos.
De acordo com os dados apresentados, todos os oito cursos avaliados com conceito 5 eram vinculados a Universidades.
Somente um curso em Universidade ficou sem conceito ( SC ).
Esse tipo de Organização Acadêmica teve o conceito 4 como modal, com 28 cursos ( 39,4 \% ).
Os demais cursos de Universidades avaliados receberam os conceitos 1 ( três cursos ), 2 ( cinco cursos ), 3 ( 26 cursos ) e conceito 5 ( oito cursos, como já mencionado ).
Somente um curso em Centro Universitário participou, e esse foi alocado no conceito 3, representando, portanto, a totalidade dessa categoria.
Nas Faculdades, apenas um dos dez cursos ficou sem conceito ( SC ), e nove cursos foram igualmente alocados nos conceitos 1, 2 e 3 ( três cursos, cada ).
Nenhum curso, neste tipo de Organização Acadêmica, recebeu conceitos 4 ou 5.
Considerando -se separadamente as regiões brasileiras, verifica -se que, na região Norte, as Universidades concentraram seis dos dez cursos participantes.
Dos cursos em Universidades, nenhum ficou sem conceito ( SC ) ou recebeu conceitos 2 ou 5.
Três Receberam o conceito modal, 3, enquanto apenas dois cursos receberam conceito 4.
E somente um curso obteve conceito 1.
Na região Norte não houve nenhum curso pertencente a Centro Universitário.
As Faculdades participaram com quatro cursos na região Norte, e, nessa Organização Acadêmica e região, somente um curso ficou sem conceito ( SC ), um curso recebeu conceito 1, dois cursos receberam o conceito 2, enquanto nenhum curso obteve os conceitos restantes nessa categoria.
Na região Nordeste, as Universidades participaram com 22 dos 24 cursos na Área de Zootecnia da região.
Nenhum dos cursos em Universidades no Nordeste ficou sem conceito ( SC ).
O conceito modal foi 3, com 11 cursos.
Os demais receberam conceitos 1 ( dois cursos ), 2 ( três cursos ), 4 ( cinco cursos ) e conceito 5 ( um curso ).
Assim como na região Norte, os Centros Universitários não ofereceram nenhum curso na região Nordeste.
As Faculdades foram representadas por apenas dois cursos na região, os dois alocados no conceito 1, correspondendo, portanto à totalidade de cursos na combinação de categoria e região.
Na região Sudeste, as Universidades concentraram 17 dos 20 cursos da região.
Entre os cursos em Universidades na região, o conceito modal foi 4 com 12 cursos, e cinco cursos receberam conceito 5.
Nenhum curso ficou sem conceito ( SC ) ou recebeu os demais conceitos nessa categoria e região.
Os Centros Universitários participaram com apenas um curso na região Sudeste, e este foi alocado no conceito 3, correspondendo, portanto, à totalidade de cursos dessa Organização Acadêmica na região.
As Faculdades foram representadas por apenas dois cursos na região Sudeste, que também se concentraram no conceito 3, e que também, representou a totalidade da categoria na região.
Dos 14 cursos da região Sul, 13 eram de Universidades, para os quais o conceito modal foi 4, com seis cursos.
Nesse tipo de organização, nenhum dos cursos ficou sem conceito ( SC ) ou recebeu os conceitos 1 ou 2.
Os demais receberam os conceitos 3 ( cinco cursos ) e 5 ( dois cursos ).
Os Centros Universitários da região Sul, assim como nas regiões Norte e Nordeste, não receberam nenhum curso.
Foi somente um, o curso vinculado a Faculdade na região Sul, e este foi alocado no conceito 3.
Na região Centro-Oeste, 13 dos 14 cursos eram de Universidades.
Nesse tipo de organização, somente um curso ficou sem conceito ( SC ), enquanto nenhum curso obteve os conceitos 1 ou 5.
Já o conceito modal foi 3, com sete cursos.
Os outros cursos obtiveram os conceitos 2 ( dois cursos ) e 4 ( três cursos ). conceitos 1 ou 5.
Já o conceito modal foi 3, com sete cursos.
Os outros cursos obtiveram os conceitos 2 ( dois cursos ) e 4 ( três cursos ).
Os Centros Universitários da região Centro-Oeste não apresentaram nenhum curso que representasse essa categoria, fato também ocorrido nas regiões Norte, Nordeste e Sul.
Somente um curso de Faculdade participou na região Centro-Oeste, e este associou -se ao conceito 2, correspondendo, portanto, a totalidade dessa categoria na região Centro-Oeste.
O Gráfico 5.3 apresenta a distribuição do Conceito ENADE segundo a Organização Acadêmica da IES.
Os cursos em Universidades ( linha azul ) apresentam uma poligonal mais à direita do que as demais, denotando uma melhor distribuição de Conceitos ENADE.
No outro extremo temos os cursos em Faculdades ( linha vermelha ) com a poligonal mais à esquerda e consequentemente a distribuição de Conceitos ENADE com menores valores.
Cursos em Centros Universitários ( linha verde ) apresentam uma situação intermediária.

% matched lemmas: 
\end{textsample}
